% The certificate stage of the classification pipeline: what each certificate proves, the thirty-fourth open
% quotient by one component group, the curves that #128 modelled proved from their branch divisors alone,
% the Gross-Zagier step of 6_323 made exact, and two corrections to the counting of CM points.
% Code: HyperellipticCertificates.m, DualTest.m, ComponentGroups.m; data: data/certificates/, data/dual/.
% Tests: tests/HyperellipticCertificates.m, tests/CMPointCounts.m, tests/OpenThirtyFour.m and
% tests/_offline/{DualTestReproduces,HyperellipticCertificatesC,HyperellipticCertificatesC4}.m.
\documentclass[11pt,reqno]{amsart}
\usepackage[margin=1.15in]{geometry}
\usepackage{amsmath,amssymb,amsthm,mathtools}
\usepackage[colorlinks=true,linkcolor=blue,citecolor=blue,urlcolor=blue]{hyperref}
\theoremstyle{plain}
\newtheorem{theorem}{Theorem}[section]
\newtheorem{proposition}[theorem]{Proposition}
\newtheorem{lemma}[theorem]{Lemma}
\theoremstyle{definition}
\newtheorem{remark}[theorem]{Remark}
\DeclareMathOperator{\Jac}{Jac}
\DeclareMathOperator{\Pic}{Pic}
\DeclareMathOperator{\Cl}{Cl}
\newcommand{\Q}{\mathbb{Q}}
\newcommand{\Z}{\mathbb{Z}}
\newcommand{\F}{\mathbb{F}}
\begin{document}
\title{Certificates for the open quotients}
\author{}
\date{\today}
\maketitle

\section{What the pipeline now proves}

The filters of the classification leave $34$ Atkin--Lehner quotients $Y=X_0^D(N)/W$ of genus at least three
undecided (\texttt{data/open\_34.m}). The last stage of the pipeline decides a curve from a \emph{certificate}: a
file of \texttt{data/certificates/} that names the inputs of a proof, every hypothesis of which the stage
rechecks. The searches that found the inputs are not repeated. Each certificate has one of the following
shapes; the first two are those of \texttt{paper/torsion-counting-hyperelliptic.tex} and
\texttt{paper/hyperelliptic-candidates.tex} (Proposition~prop:normalform), the third and fourth those of
\texttt{paper/hyperelliptic-15-38.tex} and \texttt{paper/hyperelliptic-6-323.tex}.

\begin{itemize}
\item[A] Genus three, $W$ the whole group: effective divisors of degree two whose differences are torsion
(Borcherds forms, or the dual test against saved cusp forms of weight $\frac32$), more of them than the torsion
bound of the Jacobian, lowered where needed by a component group in which the differences vanish.
\item[B] Genus four over a genus-two star: the branch divisor, recomputed from Ogg's fixed points, is linearly
equivalent to a distinct effective divisor of degree two, and the star's Jacobian is torsion-free.
\item[C] Genus three over a genus-one star $E$: the branch divisor is $A+B$ with $A-B$ torsion (the dual test,
on the $a_q(E)$-eigenspace where needed), zero (its class in the component groups $\Phi_p$, and every nonzero
rational torsion point of $E$ off the identity component at one of those $p$), and the cover untwisted (a
simple Prym, or the dual graph of $Y$). $E$ is identified exactly by its traces of Frobenius up to the Sturm
bound and by its Tamagawa numbers against the orders of the $\Phi_p$.
\item[C4] As C, with one branch place of degree four, the torsion statement from Gross--Zagier (Section~4).
\item[N] Not hyperelliptic (Section~2).
\end{itemize}

At the time of writing $26$ of the $34$ are decided this way: all seven of shape A, five of the six of shape B,
the four of shapes C and C4 that \texttt{paper/hyperelliptic-candidates.tex} lists separately, nine of the
sixteen that \texttt{paper/heegner-height-models.tex} modelled (Section~3), and the thirty-fourth (Section~2).

\section{The thirty-fourth quotient is not hyperelliptic}

\begin{theorem}
$Y=X_0^{141}(2)/\{1,6,94,141\}$ is not hyperelliptic.
\end{theorem}

\begin{proof}
$Y$ has genus three and is a double cover of the genus-one star $E=X_0^{141}(2)/W$, branched at the fixed
points of the involutions outside $\{1,6,94,141\}$: the rational CM points $P_{-4}$, $P_{-12}$ and the place
$Z_{-1128}$ of degree two. If $Y$ were hyperelliptic, its hyperelliptic involution $\iota$, defined over $\Q$
and central in $\mathrm{Aut}(Y)$, would commute with the deck involution and descend to an involution
$x\mapsto c-x$ of $E$; it would pair the four branch points into
two pairs of equal sum in the group $E$, Galois-stably (\texttt{paper/hyperelliptic-15-38.tex}, Section~1). The
only Galois-stable pairing is $\{P_{-4},P_{-12}\}\mid Z_{-1128}$, so $\Delta=P_{-4}+P_{-12}-Z_{-1128}$ would be
zero in $\Pic^0(E)(\Q)$. At $p=3$, which divides $D$ and is not split in $\Q(i)$, $\Q(\sqrt{-3})$ or
$\Q(\sqrt{-282})$, the special fibre of $E$ is a quotient of the \v{C}erednik--Drinfeld fibre, its component
group is $\Phi_3\cong\Z/7$, and the CM points specialise to the components and singular points whose Eichler
orders admit the corresponding optimal embeddings (Molina, Theorems 5.3, 5.4, 5.8). The class of $\Delta$ in
$\Phi_3$ is $2\neq0$, so $\Delta\neq0$.
\end{proof}

This replaces the argument of \texttt{paper/heegner-height-models.tex}, which needed the enumeration of the
candidate branch divisors to be complete. (\texttt{data/certificates/141\_2.m}, shape N.)

\section{The curves modelled by Heegner heights}

Each of the sixteen curves of \texttt{paper/heegner-height-models.tex}, Table 1, is a genus-three double cover of
an elliptic quotient, and when that quotient is the star, shape C proves it hyperelliptic from its branch
divisor alone, with none of the inputs that made the identification of the model conditional (the
completeness of the enumeration, the recognition of integers from $L'$-values, the Mordell--Weil groups of the
isogeny class, the oldform multiplier). Done so far: $39_{11}$, $6_{185}$, $141_2$ ($\{1,3,94,282\}$), $21_{34}$,
$55_7$, $57_{10}$, $15_{46}$, $6_{253}$, $6_{247}$. At $21_{34}$ and $15_{46}$ the four branch points are rational
and exactly one of the three pairings is certified; the other two are refused. Open: $6_{301}$, $6_{455}$ and
$690_7$ (bases being computed), $22_{51}$ (the pairing is certified zero, but neither twist test applies),
$6_{187}$ (the Galois-stable pairing is a mixed one over $\Q(\sqrt{-2})$), and $1326_1$, $210_{13}$, where $W$ has
index four and the elliptic quotient is not the star. The models of the note remain as equations; they are
theorems once the curve is proved hyperelliptic and the model is identified, which still needs the note's
enumeration or another argument.

\section{The Gross--Zagier step of $6_{323}$, exactly}

For $Y=X_0^6(323)/\langle w_6,w_{19},w_{34}\rangle$ over $E=1938h1$, the pairing $\Delta$ of the four points of
discriminant $-7752$ is the CM cycle twisted by the genus character $\chi=(-8,969)$, and it is torsion by the
Gross--Zagier formula of Yuan--Zhang--Zhang once $L(E/\Q(\sqrt{-1938}),\chi,s)=L(E^{(-8)},s)L(E^{(969)},s)$
vanishes to order three at $s=1$. $E^{(-8)}$ has root number $-1$. $E^{(969)}$ has root number $+1$, and
$L(E^{(969)},1)=0$ is checked exactly: for the newform $f$ of $E$ and the primitive quadratic character
$\varepsilon$ of conductor $m=969$,
\[
  L(f\otimes\varepsilon,1)=\frac{\tau(\varepsilon)}{m}\sum_{a\bmod m}\bar\varepsilon(a)\,
  2\pi i\!\int_{a/m}^{i\infty}\! f(z)\,dz
\]
(Birch), which needs no coprimality between $m$ and the level, and $f\otimes\varepsilon$ is the $L$-series of
$E^{(969)}$: at the primes $3,17,19$ of $m$, where $E$ is multiplicative, $\varepsilon(p)=0$ and the twisted curve
has $a_p=0$. So $L(E^{(969)},1)=0$ exactly when the image of the twisted winding element
$\sum_a\varepsilon(a)\{0,a/m\}$ in the $f$-part of the modular symbols of level $1938$ is zero, which it is. The
same computation gives nonzero values for the twists of root number $+1$ by $5$, $17$, $21$. This replaces the
proved rank $2$ of $E^{(969)}$ that the earlier argument used, which does not by itself give the vanishing.

\section{Two corrections to the counting of CM points}

\emph{The Heegner weight of $P_{-4}$.} The dictionary of the dual test weights a CM point $P_d$ by $1/e_d$, $e_d$
the order of its stabiliser. The points of discriminant $-4$ are fixed by $w_2$ when $2\parallel DN$, since
$1+i$ has norm $2$ (Ogg), so $e_{-4}=4$ there; the earlier scripts doubled $e_d$ only for $d=-4m$ and for $d=-m$
with $m\equiv3\pmod 4$. The weight enters only targets that contain $P_{-4}$, which no earlier proof had; three
of the curves of Section~3 do.

\emph{Points over a star point.} The number of points of $X_0^D(N)$ over a point of $X_0^D(N)/W$ is $|W|$ divided
by the order of the stabiliser of the point in $W$; the units of the CM order fix the point and do not enter.
The component-group code had folded them in, which counted $P_{-3}$ on $X_0^6(247)^*$ as four points against its
degree one, and was right at $P_{-4}$ of the thirty-fourth quotient only because the two errors cancelled there.
\texttt{tests/CMPointCounts.m} checks both against Ogg's fixed-point counts and against the degrees of the fields
of definition.

\end{document}
